\chapter{Design of Embedded Systems}
In this chapter, we go over different practices and technologies regarding embedded systems.
Most notably, sensors, actuators, embedded processors, memory architectures, input output (IO), scheduling, and a bit of Lingua Franka---an embedded systems programming language.

\section{Sensors and Actuators}
Sensors are devices which measure a physical quantity, while actuators modify a physical quantity.
One vitally important sensor is the \emph{accelerometer}, which measures acceleration, but it can also be used to measure the tilt of an embedded system.
To achieve this, optical gyroscopes can also be used.
Accelerometers can also be used to track a rough estimation of position based on measured accelerations, this is called an \emph{intertial navigation system} (INS).

A problem all sensors have is, that they are notoriously known for \emph{drift}, which makes calibration vitally important.
To analyse this problem and other inaccuracies of a sensor, we define the \emph{affine sensor model}.

\begin{definition}[Affine Sensor Model]
    The affine sensor model defines physical properties as values \(x \in \R\) and sensors as so-called \emph{affine} linear functions \(F : \R \rightarrow \R, x \mapsto ax + b\), with \(a, b \in \R\), where \(a\) is called \emph{sensitivity} and \(b\) is called \emph{bias}.

    To further enhance the realism of our model, we consider a \emph{saturated} range \(L, H \in \R, L \le H\), and define the sensor function as
    \begin{IEEEeqnarray*}{c}
        F : \R \rightarrow \R, x \mapsto
        \begin{cases}
            ax + b & \text{if} x \in [L, H] \\
            aH + b & \text{if} x > H \\
            aL + b & \text{if} x < L
        \end{cases}.
    \end{IEEEeqnarray*}
    As a last enhancement, we consider the \emph{precision} \(p \in \R\) of the sensor, which is the smallest absolute difference between sensor readings, thus an even more realistic function is \(F'(x) = p \cdot \left\lfloor F(x) / p\right\rfloor\), which clamps everything into a partitioned interval.

    The dynamic range \(D := \frac{H-L}{p}\) is often notated in decibels \(D_\textit{dB} := 20\log_{10} D\).
    A quantized signal like \(F'\) often is an \(n \in \N\) bit signal, which results in \(D_\textit{dB} := 20n \log_{10} 2 \approx 6n\).
\end{definition}
When using sensors, it is advised to carefully use them if they are prune to aliasing.
Alternatively, use higher-precision sensors.

Another way to handle sensor inprecision is to use multiple different sensors for the same quantity, which is called \emph{redundancy}.

A classic way to drive actuators is pulse-width modulation (PWM), which uses repeated small pulses of the same voltage.
Increasing pulse frequency, increases delivered power to the actuators.

\section{Memory Architectures}
Memory is used for different tasks in embedded systems, e.g., storage or communication with sensors and actuators.
Embedded systems memory is often more restrictive than general purpose memory in size, power, and reliability.
There are different types of memory.
\begin{description}
    \item[Non-volatile] preserves contents when power is off \hfill
        \begin{description}
            \item[Flash] Electrically erasable programmable read-only memory.
                Erases a block at a time.
                There is a limited number of write/erase cycles, i.e., around 100000.

            \item[Disks] Too slow for embedded systems.
        \end{description}
    \item[Volatile] loses contents when power is off \hfill
        \begin{description}
            \item[Registers] Fastest access time possible, extremely expensive and power hungry.
                Used directly inside processors.

            \item[SRAM] Static random-access memory.
                Very fast, deterministic access time.
                Expensive and power hungry.
                Used for (processor) caches and scratchpads.

            \item[DRAM] Dynamic random-access memory.
                Slower than SRAM, but less power hungry and expensive.
        \end{description}
\end{description}

Typically, embedded systems chips use memory-maps to partition all memory and sensor/actuator interfaces into different address spaces.
E.g., using \lstinline|0x0000| to \lstinline|0x1FFF| for flash, \lstinline|0x2000| to \lstinline|0x4FFF| for DRAM, and \lstinline|0x5000| to \lstinline|0x5FFF| for SRAM.

Sometimes controllers use an e.g., 8-bit architecture, but 16-bit addresses.
In this case, the controller needs two cycles to determine the position in memory for an address.

It is possible to manage the program memory in different methods.
C programs typically use three different allocation methods for this.
\begin{description}
    \item[Static Memory] The compiler chooses the address at which to store a variable.
        A variable is static, if it is a global variable, or when it is declared \lstinline|static| inside the scope of a function.

    \item[Stack] The program dynamically allocates the variable using a LIFO strategy, pushing new allocations when entering a scope and popping on leaving a scope.
        Thus in C, stack-allocated variables are any non-static variables in a function scope.

    \item[Heap] The developer manually dynamically allocates and deallocates memory on an extending structure.
        This is done by allocating address spaces with \lstinline|void* malloc(size_t size)|, and deallocating it with \lstinline|void free(void* ptr)|.
\end{description}
Sometimes stack and heap grow in the direction of the other, which limits (process) memory, but on embedded systems, a program typically can use all memory, as neither operating system nor processes exist.

Another topic to note is memory hierarchy.
Most systems implement this using a linear hierarchy: Flash can be read into memory, memory caches parts of page lines, and registers load individual values.
Scratchpads are developer-managed and typically used like cache.
